% Generated draft: requires mathematical and source review.
\documentclass[11pt]{article}
\usepackage[margin=25mm]{geometry}
\usepackage{fontspec}
\setmainfont{Latin Modern Roman}
\usepackage{amsmath,amssymb,mathtools,unicode-math}
\setmathfont{Latin Modern Math}
\usepackage{enumitem,needspace,xurl,hyperref}
\setlength{\parindent}{0pt}
\setlength{\parskip}{5pt}
\setlength{\emergencystretch}{3em}
\newcommand{\N}{\mathbb{N}}
\newcommand{\Z}{\mathbb{Z}}
\newcommand{\Q}{\mathbb{Q}}
\newcommand{\R}{\mathbb{R}}
\newcommand{\C}{\mathbb{C}}
\newcommand{\F}{\mathbb{F}}
\newcommand{\A}{\mathbb{A}}
\newcommand{\PP}{\mathbb{P}}
\newcommand{\E}{\mathbb{E}}
\newcommand{\eps}{\varepsilon}
\newcommand{\dd}{\,\mathrm{d}}
\begin{document}
\section*{UnsolvedMath ID: 11000103}
\textbf{Source problem code:} catalogue entry
\subsection{Definitions and mathematical statement}
Let $S$ be the surface in the source's mapping-class-group setting and let $\operatorname{Mod}_S$ be its mapping class group. For $f\in\operatorname{Mod}_S$, write
\[
 \langle\!\langle f\rangle\!\rangle=\langle hfh^{-1}:h\in\operatorname{Mod}_S\rangle
\]
for its normal closure. A pseudo-Anosov mapping class has stretch factor $\lambda(f)>1$.

\textbf{Conjecture I (source chapter of N. Ivanov).} A pseudo-Anosov $f$ with sufficiently large dilatation should have normal closure a free group with a free generating set consisting of conjugates of $f$. A natural surface-uniform reading of ``sufficiently large'' is that there is a threshold $\Lambda(S)$ such that the assertion holds whenever $\lambda(f)>\Lambda(S)$; the excerpt does not specify the permitted threshold dependence or the class of surfaces.

The source offers a more cautious power version: for a given pseudo-Anosov $f$, a sufficiently high power $g=f^N$ should have a free normal closure with a free basis drawn from conjugates of $g$. Its basic existence formulation is
\[
 \forall f\text{ pseudo-Anosov}\quad\exists N\ge1\quad
 \langle\!\langle f^N\rangle\!\rangle\cong F(B),
 \quad B\subset\{hf^Nh^{-1}:h\in\operatorname{Mod}_S\}.
\]
Here the isomorphism sends each named free generator to the corresponding conjugate. Whether ``sufficiently high'' is intended to assert all $N\ge N_0(f)$, or a specified divisibility condition, is not stated by the excerpt.

The phrase ``having as generators the conjugates'' needs care. All conjugates generate the normal closure by definition; the substantive free-basis conclusion concerns a suitable selected family. The entire conjugacy class cannot silently be treated as a free basis, since conjugation by elements of the normal closure produces relations among that whole family. The source does not supply a transversal or indexing rule for a basis.

The source also suggests that every nonidentity element of such a subgroup might be pseudo-Anosov. This is an additional purely-pseudo-Anosov expectation, not a consequence of freeness alone. It motivates the conjecture by Gromov [Gr], Theorem 5.3.E, on normal closures of hyperbolic elements, and Delzant [D] for a more direct approach. The proposed Teichm\"uller-space negative-curvature-along-axes analogy is motivational, not a specified curvature hypothesis or theorem. The reference to ``Problem 3 above'' has no reproduced target here.
\subsection{Short English statement}
Ask whether sufficiently large-dilatation pseudo-Anosovs, or at least sufficiently high powers of each one, normally generate free groups with bases made of selected conjugates. The source additionally hopes these subgroups are purely pseudo-Anosov.
\subsection{Sources}
\begin{itemize}
\item Local UnsolvedMath registry, permanent ID 11000103, source code catalogue entry; no separate proposer or proposal year supplied.
\item Farb, edited volume \emph{Problems on Mapping Class Groups and Related Topics}, Conjecture I, PDF page 86, in source text bearing N. Ivanov's chapter attribution; \url{https://www.math.uchicago.edu/~farb/papers/mcgbook.pdf}. The registry records PDF-font-structure extraction and access on 2026-07-29; the inventory distinguishes 329 labeled headings by source page.
\item M. Gromov, source label [Gr], Theorem 5.3.E, and Th. Delzant, source label [D]. Full bibliography entries and exact hypotheses are absent from the record.
\end{itemize}
\end{document}
